





基于高保真模拟器的车辆与飞行器的协同控制
\section{绪论}

\subsection{研究背景与意义}

在一次航班离港流程中，飞机从停机位推出到进入滑行路线，看起来只是
周转链条中的一个短环节，但它实际上同时牵涉机坪交通、廊桥区域安全、地
勤人员协同和后续放行节奏。牵引车需要在飞机前起落架附近低速靠近，车身
周围又存在机翼、发动机、廊桥、保障车辆和地面标线等多种约束。任何一次
偏航过大、停车不及时或目标识别错误，都可能使简单的推出作业演变为安全
事件。因此，飞机牵引与推出并不是单纯的“车辆驾驶”问题，而是一个典型
的低速、高精度、多约束协同作业问题[1-2]。
目前多数机场的推出与牵引仍以人工驾驶牵引车、机务或指挥员目视协同
为主。该方式在经验丰富、交通密度较低的情况下较为可靠，但其稳定性很大
程度取决于人员状态和现场沟通质量。航班高峰、夜航、雨雾天气以及机坪车
辆交叉运行时，人工判断容易受到视距、噪声和遮挡影响，牵引路径规划、翼
尖安全距离判断、前起落架对准等环节都可能出现误差。对智慧机场而言，这
类高度依赖人工经验的作业环节，会限制地面运行系统继续向闭环调度和自动
执行方向发展。
相关资料表明，机场地面运行阶段的安全风险并不低，推出、牵引和低速
滑行是典型的风险集中环节之一，常见问题包括牵引车与机体接触、翼尖剐
蹭、人员误入作业区域等。同时，地面运行延误还会沿航班计划向后传递，使
单个机位或单个航班的问题进一步影响后续起飞排序和机坪资源利用。因此，
提高牵引作业的自动化水平，不仅有助于降低局部操作风险，也有助于提升机
场地面运行的整体稳定性[3-6]。
从机场运行系统的发展方向看，SMGCS 向 A-SMGCS 升级后，场面目标监
视、路线引导和冲突提示逐渐被纳入统一框架[2-4]。这些研究能够回答“机坪
上有哪些目标、目标应如何运行、是否存在冲突”等问题，但本文关注的位置
更靠近执行端：牵引车已经进入飞机前起落架附近时，如何控制转向、速度和
停车距离。也就是说，上层 A-SMGCS 可以给出任务和安全边界，而末端对接仍
需要车辆级建模、感知和控制方法支撑。






基于上述背景，本文选择以高保真仿真环境为研究支撑，围绕机场牵引车
与飞行器的协同控制展开建模、感知、路径规划和控制验证。通过在 ROS 与
Gazebo 环境中构建牵引车、飞机目标和机坪作业场景，可以在较低成本下复现
实车试验中难以频繁开展的对接过程，观察车辆转向、速度、目标识别和安全
约束之间的相互影响[11-13]。这对于后续自动牵引系统的样机设计、控制算法调
试以及机场地面作业智能化研究都具有现实意义。
\subsection{国内外研究现状}
\subsubsection{机场牵引与地面自动化研究现状}
机场地面自动化研究的出发点，是在不降低安全裕度的前提下减少人工重
复操作，并提高机坪资源的利用效率[1-4]。飞机推出与牵引处在航班周转和场面
交通之间：一方面，它需要服从停机位、滑行通道、速度限制和指挥流程等机
场运行规则；另一方面，它又要求牵引车在飞机前起落架附近完成低速对准、
连接和稳定牵引。因此，该环节既有车辆控制问题，也有机场运行组织问题。
国外研究较早把飞机地面滑行阶段作为节能和自动化对象。TaxiBot、e-
Taxi 等方案的共同思路，是尽量减少飞机发动机在地面工作时间，从而降低油
耗、噪声和排放[6,8-9]。Zaninotto 等则从牵引车调度与路线安排角度，利用
冲突检测和时序优化减少车辆之间、车辆与飞机之间的干扰[7]。另一些工作关
注牵引车—飞机组合体的制动、转向和纵向响应，为低速牵引控制提供了动力
学参考[6,10,19]。
从机场系统层面看，SESAR 以及 A-SMGCS 相关研究主要解决的是“场面上
有哪些目标、目标往哪里走、是否存在冲突”这类上层运行问题。监视、路径
引导、冲突预警和运行管理被放在统一框架内后，机场能够更细致地掌握航空
器和地面车辆状态[2-4]。不过，这类框架通常只给出任务、路径和安全边界，
并不会直接处理牵引车贴近前起落架时的转向角、停车距离和连接姿态。因
此，它能够为自动牵引提供运行环境，却不能替代末端对接阶段的车辆级控制
研究。
国内在飞机牵引滑行和地面运行智能化方面起步较晚，但近几年已出现较
多综述和工程研究。孙艳坤等梳理了牵引装备、动力学特性和运行模式，并将
自动牵引与智能调度列为后续发展方向[6]。中国民航局发布的 A-SMGCS 相关指
南和标准，也为场面运行信息化提供了制度依据[1-2]。从公开文献看，现有成
果多集中在系统规划、装备类型或单项算法讨论，对“牵引车靠近前起落架并







停在合适位置”这一末端过程的仿真验证仍显不足。
基于此，本文没有把研究范围铺得很大，而是把问题收束到一个具体作业
单元：在给定飞机位置、车辆结构和机坪规则后，牵引车如何确认目标、生成
可执行接近路线，并在低速状态下完成对准和停车。该问题与实际调试更贴
近，也便于在 ROS-Gazebo 中反复改变初始位姿、控制参数和传感器布置进行验
证。
\subsubsection{协同控制与路径规划方法研究现状}
从方法上看，机场牵引自动化可归入移动机器人导航和智能车辆控制范
[17-20,28]
畴，但它又不能完全照搬普通道路自动驾驶或仓储 AGV 的方案                。原因在
于，牵引车面对的是航空器这一高价值目标，运行区域空间狭窄、规则刚性
强，而且最终对接阶段对横向误差和姿态误差都十分敏感。也就是说，路径规
划不仅要“走得到”，还要“转得过、停得住、对得准”。
图搜索与采样类算法是路径规划研究中的基础方法。A*、Dijkstra 等图搜
索算法结构清晰、结果可解释性强，适合在机场滑行道、机坪通道等规则环境
中生成全局参考路径[22-23]；RRT 及其改进算法则更适合在连续空间中快速搜索可
行路径，常用于机器人和自动车辆的运动规划[20-21,24]。这类方法的优势在于稳定
和易于检查，但在动态障碍较多、局部空间受限或车辆曲率约束较强时，单独
使用往往需要额外的平滑、重规划和可行性筛选。
智能优化算法也可用于路径规划，例如把距离、安全间隔、能耗或时间写
入同一个目标函数，再通过迭代搜索较优路线。不过，这类方法对权重、初始
种群和停止条件较敏感。对于机场牵引这种安全边界很明确的任务，如果同一
场景多次求解得到的轨迹差别较大，或者规划依据难以复查，后续工程使用就
会比较受限。
强化学习和深度强化学习在动态决策中具有吸引力，能够通过仿真交互学
习避障或转向策略。但机场牵引不是单纯追求“能跑起来”的场景，安全可验
证性和监管可解释性更重要。若完全依赖黑箱策略，训练样本、仿真真实性和
迁移误差都会成为风险点；一旦仿真机坪与真实机坪存在差异，学到的策略未
必能直接用于实车。
因此，本文更倾向于采用分层融合方案。上层保留机坪规则、曲率上限和
安全距离等显式条件，使参考路径能够被检查；下层再根据传感器给出的纵
向、横向和航向误差修正车速与转角。这样的处理不依赖单一智能算法，也避







免纯规则方法在近距对接阶段反应不足，为牵引车—飞机协同控制留下了较稳
妥的实现路径[25,28]。
\subsection{研究内容与技术路线}
本文以机场飞机推出与牵引作业为应用背景，重点研究牵引车与飞行器在
低速近距场景下的协同控制问题。课题并不追求单一算法指标，而是围绕“建
模—感知—规划—控制—仿真验证”这条链路展开：先把机场牵引作业抽象成
可计算的问题，再利用仿真平台检验车辆能否按预期完成目标识别、接近、对
准和停车。结合本文已经开展的 ROS/Gazebo 建模、SolidWorks 车辆模型导
入、双桥转向控制调试等工作，主要研究内容如下。
\paragraph{} （1）机场牵引作业流程建模与问题抽象
从停机位推出和前起落架对接过程出发，分析牵引车、飞机、机坪通道、
禁行区域和障碍物之间的关系。根据作业顺序，将自动牵引过程划分为目标确
认、初始接近、姿态对准、微动对接和停车保持等阶段，并把实际约束转化为
状态变量、控制输入和安全边界。该部分的目的，是为后续控制算法提供一个
统一的问题描述，而不是只停留在场景文字说明上。
\paragraph{} （2）面向自动对接的飞机辅助感知方法
飞机辅助感知部分主要解决目标确认和误差输入问题。本文将前起落架附
近区域作为识别对象，在仿真中设置辅助标识，使车辆能够读出目标身份、相
对距离、横向偏差和姿态偏差；视觉模块负责识别和方向估计，激光雷达、超
声或等效近距传感器用于校核安全距离。该部分重点回答三个工程问题：接近
哪架飞机、从哪个方向对准、何时必须减速或停车。
\paragraph{} （3）多约束条件下的路径规划与协同控制方法
控制方法部分以第二章运动学模型和感知误差为输入。全局层面根据机坪
通道、车辆曲率和安全距离生成可执行参考目标；局部层面根据实时误差调节
速度和转向。控制设计同时考虑低速稳定性、转角限幅、障碍距离、终端停车
阈值和姿态对准要求。本文将模式管理、局部目标、误差反馈和执行器映射分
开实现，便于后续在 ROS 节点中调试。




\paragraph{} （4）基于高保真模拟器的系统集成与验证
依托 ROS 与 Gazebo 搭建牵引车—飞机对接仿真系统。车辆模型由








SolidWorks 建立并经 URDF/STL 等形式导入 Gazebo，前后桥转向关节和四个驱
动车轮通过控制器实现位置和速度控制，传感器模型挂载于车体前部。通过设
置典型机场作业场景，对牵引车模型加载、转向回中、轮速控制、目标接近和
停车判定等功能进行验证，并从安全性、路径平滑性和控制稳定性等角度分析
系统表现。
\paragraph{} （5）牵引车辆结构与动力系统的工程化思考
在完成仿真和控制算法研究的同时，结合机场牵引车实际需求，对车辆结
构、动力布置和维护便利性进行初步分析。该部分不作为完整机械设计展开，
而是从工程实现角度讨论车辆硬件结构与控制策略之间的关系，例如双桥转向
结构对小半径机动和精细对位的影响、动力响应对低速控制稳定性的影响等，
为后续实物样机或半实物平台研究提供参考。


